\chapter{总结与展望}\label{chap:conclusion}

\section{工作总结}

本课程设计以「人机交互导论」课程要求为出发点，围绕「让音乐推荐理解用户的此刻情绪」这一目标，完成了 Mood Radio 这一面向本地音乐库的情绪感知电台 Web 应用的设计与实现。回顾这一过程，工作覆盖了从设计研究、系统架构、关键算法到交互细节的完整链路。

在设计层面，课题组并未一上来就动手写代码，而是先用问卷、情景访谈、焦点小组与单独访谈四种方法做了一轮规模虽小、但相互印证的用户研究，从 42 名调研对象的反馈中提炼出三个 Persona，以及「低负担输入、用户可控、可解释、本地优先、沉浸感与一致性」这五条贯穿全局的原则。这一前置工作让后续的功能取舍始终有依据可查，避免了「先做出来再说」的常见陷阱。在工程层面，系统在三周左右的时间内完成了 React 前端的 8 个页面、25 个以上 REST 接口的 C++ 后端、SQLite 持久层与 Python 离线分析脚本的协同实现；其中包括了对 ncnn 推理失真这一关键工程问题的诊断与绕过——这是一段相当耗精力的过程，但它让最终的曲库标注重新可用，也成为本项目工程成熟度的一个侧面证明。在交互层面，五种情绪输入通道（摄像头表情、手动选择、自然语言聊天、时段电台、头部手势）配合 7 套情绪主题、Aurora 流动背景与黑胶旋转封面，构建了一个用户可以\emph{感知}到的反馈环；课堂模拟评估显示，9/10 的受试者认可本地隐私设计，8/10 的受试者能够理解推荐逻辑——这两个数字并不意味着系统已经完美，但它们至少证明了课题组的关键设计判断是站得住脚的。

整体上看，本项目可以被看作是把人机交互课程中的「多模态输入、情境感知、用户控制、可解释性、隐私保护」这些抽象概念，翻译成一份具体的、可被使用与检验的产品形态的一次尝试。如果要为这次尝试找一句概括，那么大概是：「\emph{人机交互的真正难点不在于识别用户的情绪，而在于识别之后如何尊重它}」。

\section{局限性}

诚实地看待这一作品，仍有不少明显的局限。最为底层的一项是 \textbf{C++ 端的 ncnn 推理路径仍然不可用}——目前所有 V/A 标注都依赖 Python 离线脚本，这虽然在课堂演示中工作良好，但它意味着「新上传文件的实时分析」必须通过子进程来绕一圈，部署也比理想方案多一层依赖。要彻底解决这一问题，需要重做 ONNX 到 ncnn 的转换并把 FeatureExtractor 与训练流水线严格对齐，这是后续值得投入的工程工作。

\textbf{群体在场感不够实时}是第二项局限。系统目前通过 60 秒心跳上报与 10 秒前端拉取实现群体页面更新，端到端的延迟接近一分钟；这一节奏对「我想看看现在多少人在听同一种情绪的歌」这种弱时效性需求是够用的，但如果想做到「看到别人开始播这首歌的瞬间高亮闪烁」这一类更接近真实多人协同的体验，就必须引入 WebSocket。课题组在课程时间内没有为此投入精力，是经过权衡的选择，但也意味着这一方向仍有不少可挖掘的空间。

\textbf{推荐画像偏简单}是第三项局限。当前的复合打分对「用户喜欢」只加固定的 1.5 分，并没有基于用户长期反馈动态地调整目标 V/A 中心。换句话说，系统目前能记住「这首我喜欢」，但不知道「我喜欢的歌大致集中在哪一带」。要让推荐随时间真正个性化，需要在打分中引入「用户偏好画像」这一项，并允许它在每次反馈后做小步更新。

\textbf{评估规模有限}是第四项局限。$N = 10$ 的课堂模拟评估只能提供方向性结论，无法给出严格意义上的统计显著性。若要进一步检验设计判断，需要扩展为正式的 within-subject 实验，对推荐准确率、操作时间、主观满意度等指标做更细致的对比分析。

最后，\textbf{弱光识别与小曲库问题}是两项分别在两个维度上的具体不足：face-api.js 在弱光、遮挡场景下识别不稳定，而当个人曲库较小时，复合打分容易反复推送少数高分曲目。两项的解决思路不同——前者属于算法层面，后者属于多样性 / 探索的产品策略层面——但它们都已经在评估章节中得到了用户反馈的印证。

\section{后续可扩展工作}

围绕上述局限，本项目在后续可以沿着若干方向继续推进。最为根本的方向是\textbf{修复 ncnn 推理}，让 C++ 路径真正可用，从而摆脱对 Python 子进程的依赖，简化部署。\textbf{偏好画像 v3} 则是推荐侧的下一步：把所有「喜欢」歌曲的 V/A 均值视作「用户喜欢的中心」，作为额外的打分项参与复合排序，并随用户反馈动态更新。如果有更多时间，引入 \textbf{WebSocket 在场推送}可以把「群体在场」从「每分钟刷新一次」升级到「实时同步」，给「正在听同一首歌」这种弱社交体验提供更高的呈现密度。\textbf{多设备同步}也是一个相对成熟的方向——通过 UUID 关联，跨浏览器同步播放历史与喜好。

更具想象空间的几个方向则与「沉浸感」有关。\textbf{节拍/BPM 同步}让黑胶旋转的速度跟随歌曲 BPM——一首 120 BPM 的歌的封面会比一首 60 BPM 的封面转得快两倍——可以进一步把视觉沉浸推向感官细节层面。最后，\textbf{推荐解释自动生成}是另一个值得探索的方向：让 LLM 在每次推荐结果旁边实时生成一段「为什么推这首」的解释短文。这一能力如果做得好，会让本项目的「可解释性」从「徽标」这一冷静的形态，升级到「自然语言陪伴」这一温暖的形态——而这正是「机器理解情绪之后如何回应」的一个自然推进方向。

\section{结语}

Mood Radio 在课程设计的尺度内做了一次完整的「从识别到回应」的人机交互探索。如果说我们从这次探索中收获了一条最值得记下来的判断，那就是：好的情绪交互系统不应只停留在「\emph{识别情绪}」，更重要的是\emph{识别之后如何回应、如何解释，并始终尊重用户的控制权与隐私边界}。识别本身越是精准，这一立场就越关键——因为系统越懂用户，用户对系统的边界感就越敏感。本项目希望传达的正是这种立场。
